\subsection*{先找两个根的中心}
二维特征多项式是 $\lambda^2-(a+d)\lambda+(ad-bc)$。不急着套二次公式，先把未知数移到两根的平均位置。令 $m=(a+d)/2$，写 $\lambda=m+r$，一次项便抵消，剩下的只是偏移距离 $r$。

整理得到 $r^2=((a-d)/2)^2+bc$。对角元差和交叉作用的乘积共同决定偏移；迹的一半决定中心。配方因此揭示了结构，而不只是缩短算式。平方偏移的正负还告诉我们实根存在与否。

\subsection*{快速求根之后检查什么}
对 $\left[\begin{smallmatrix}4&2\\1&3\end{smallmatrix}\right]$，中心是 $3.5$，平方偏移为 $2.25$，根是 $5,2$。检查它们的和等于迹 $7$，积等于行列式 $10$。两项核对能发现不少算术或符号错误。

若根号内为负，在实数域没有根，在复数域仍有一对共轭根。旋转动作完全可以很自然，却没有实的保留方向；这个例子提醒我们区分“实数中无解”和“任何意义上都无解”。

若根重合，还必须看特征空间维数。单位矩阵保留所有方向，而非平凡剪切只有一条特征直线，虽然两者都能有重根。这个方法加快的是特征值计算，没有省掉判断可对角化的工作。把和、积、特征方向三种检查一起使用，快速计算才不会掩盖问题结构。
